\section{Method details}
\label{app:method_details}
This appendix holds the parts of the method that the main text states briefly: the choices behind the pathwise estimator and the critic (Sections~\ref{sec:method} and~\ref{sec:jpc_2d}), the design of the joint training (Algorithm~\ref{alg:jpc}), and the amortization loop in full. Notation is that of the main text: $s_k=(\rvz_{t_k},\rvx_{t_{k-1}})$, the denoiser $\mathrm d_\theta$, the sampler coefficients $\lambda_k,\eta_k,\psi_k$ of~(\ref{eq:sde}), the reward $R(s_{K+1})$ of the decoded response, and $\mathcal G$ for the set of generated positions.

\subsection{The pathwise estimator and the critic}
\label{app:pw_details}
\paragraph{One transition.}
The estimator of $\mathcal J_{\rm SVG}$ draws one transition $k$ uniformly from a window $\mathcal W=\{k_{\rm lo},\dots,k_{\rm hi}\}$ of the sampler, with $1\le k_{\rm lo}$ and $k_{\rm hi}\le K-1$, and differentiates the critic through that transition alone:
\begin{equation*}
\hat g_{\rm PW}=\nabla_\theta V_\phi(s_k),\qquad k\sim\mathcal U(\mathcal W),\qquad s_k=\big(\lambda_k\rvz_{t_{k-1}}+\eta_k\,\mathrm d_\theta(s_{k-1})+\psi_k\epsilon_k,\ \mathrm d_\theta(s_{k-1})\big),
\end{equation*}
with $s_{k-1}$ and $\epsilon_k$ held fixed (eq.~\ref{eq:pwgrad}). It is computed by rolling the sampler to $s_{k-1}$ without gradients and taking the denoiser call at $s_{k-1}$ with gradients, so its memory is that of the backward passes through the actor's call (two evaluations under classifier-free guidance) and through the critic, independent of $K$: the dependence of the reward on $s_k$ through the remaining steps is carried by the value in expectation rather than by back-propagation through them (Appendix~\ref{app:pw_derivation}). The window is the middle of the trajectory, $\mathcal W=[K/8,\,5K/8]$ on every task: it excludes the first steps, where the state is noise and the value flat, and the last, where the state is committed and the reward of the decoded response is piecewise constant; its ends are hyperparameters, not derived. With a learned critic the update is biased by the difference between the critic's slope and the true value's slope on the states the current policy visits, which is why the critic is refit on those states after every actor step (Algorithm~\ref{alg:jpc}). The update has no baseline, advantage or trust region: a baseline is a function of $s_{k-1}$, which the differentiated call holds fixed, and so drops from a pathwise gradient, and the policy's drift is controlled by the step size and by the refit of the critic.

\paragraph{Training target.}
No intermediate labels are needed. States are taken from rollouts of the sampler together with the rollout's reward $R(s_{K+1})$, and the critic is trained by binary cross-entropy, (\ref{eq:critic_loss}) for the 1D critic and~(\ref{eq:critic2d_loss}) for the 2D critic; the minimiser of either over all functions of the state is the success probability $V^\star_k(s_k)=p_\theta(R(s_{K+1})=1\mid s_k)$. \emph{Pre-training.} Before RL the critic is trained on rollouts of the task-SFT policy, with states stored at several steps spanning the trajectory and each state labelled by its rollout's reward; the copy of the denoiser and the head $f_\phi$ are trained together for a fixed number of epochs, and the epoch is chosen by a held-out score of the critic (Appendix~\ref{app:details}). The rollouts are drawn under the training sampler, or, when its early states carry no information about the outcome, under a cheaper sampler of the same denoiser, whose states the critic reads at their own time. \emph{During RL} the copy is frozen and only the head is trained. This is a cost and stability trade-off: refitting the copy after every actor step would cost a full fine-tuning step and would move the actor's objective between actor steps, whereas the head is a small correction on fixed features that a few steps refit.

\paragraph{Per-position dependence is not restricted to be affine.} With the features held fixed, if a head on the mean feature could be written as a mean of per-position readings, it would be affine (Proposition~\ref{prop:additive}(a)): a head on the mean weighs the positions only linearly, through the single bottleneck $\bar h$; a head applied per position may depend on each reading through any function, for instance one that saturates when a reading is confident.

\subsection{Design choices of the joint training}
\label{app:joint_details}
Algorithm~\ref{alg:jpc} refits the head after every actor step. Its remaining choices are the following trade-offs.

\emph{Where the head is refitted.} Each iteration rolls out a fresh batch of prompts to the end of the trajectory, decodes and verifies them, and pushes their state at the top of the window, $s_{k_{\rm hi}}$, with the reward into a replay buffer; the head is then trained on minibatches from the buffer with~(\ref{eq:critic2d_loss}) (its first term for the 1D critic). The refit is made at one step only, the top of the window, where the outcome is most determined by the state, so that the head follows the policy's changes rather than the residual randomness of the outcome. The critic at the earlier steps of the window is the pre-trained one, corrected only through what is shared, the head and its time input, on features the joint phase does not change. This is a design choice for the stability of the refit, not a consequence of the objective, and its cost is that the check below tests the critic at $k_{\rm hi}$ only.

\emph{Staleness of the buffer.} The buffer holds the newest $N_{\mathcal B}$ states, pushed $B$ per iteration, so it spans $N_{\mathcal B}/B$ iterations of past policies and the head is fitted to a mixture of their values. This is a second approximation, made for the sake of the sample size of the refit; with the values of Appendix~\ref{app:details} the span is $512$ iterations, a third to a half of a run, so the head is fitted to a genuine mixture of past policies rather than to the current one. What makes it usable is that the actor's step is small and that the check below is made on the newest states only.

\emph{The critic stays ahead of the actor.} A fixed number $m_{\rm c}$ of head steps per actor step suffices while the policy moves slowly. Every $n_q$ iterations the critic's discrimination is measured on the newest $N_q$ buffer states by the AUC, the area under the ROC curve of $V_\phi$ as a score for the reward (the probability that a random successful state scores above a random failed one). If it falls below a floor $a_0$, the critic takes $m'_{\rm c}$ extra steps before the actor moves again. The AUC at $k_{\rm hi}$ is a proxy for a necessary condition, that the critic still separates outcomes under the current policy, not a test of slope agreement across the window.

\emph{Step size.} The bias of a learned critic is paid at every actor step in proportion to the step, which motivates a decreasing step; the schedule itself, a cosine decay to a tenth of the initial value, is a choice. The critic's learning rate is constant.

\emph{One sampler.} The value function is defined by the sampler: a different sampler ($\gamma$, $K$) induces different trajectories and different values, so critic and actor are trained under the sampler that will be deployed, with one exception: on Countdown the critic's \emph{pre-training} uses states of the cheap sampler, because the deployment sampler's early and middle states carry almost no information about the outcome (Appendix~\ref{app:details}); the joint refit and the actor still run on the deployment sampler. Appendix~\ref{app:loop} also uses the critic to rank and nudge candidates on a cheaper sampler; that use rests on the assumption that a state's reading, at its own time, ranks candidates on that sampler too, which the training argument does not cover.

\subsection{The amortization loop}
\label{app:loop}
The value and the critic are defined by the sampler under which the policy was trained, and the RL improvement is made along that sampler's trajectories; a cheaper sampler, with fewer steps or $\gamma=0$, runs different trajectories through the same denoiser, and an improvement along one need not carry over to the other. The loop moves the improvement into the model's supervised objective, which trains the denoiser at every time $t$ rather than along one sampler, and uses the critic of \methodSVGTwoD{} a second time, to steer generation before the verifier can.

\paragraph{One round.}
Let $p_\theta(\rvw\mid\rvc)$ be the distribution of decoded responses under the sampler that draws the round's samples. A round takes a pool $\mathcal P$ of prompts without labels, disjoint from the evaluation prompts, draws $n$ responses per prompt, keeps the verified ones, and fine-tunes the model on them:
\begin{align}
\mathcal A(\theta)&=\Big\{(\rvc,\rvw^{(j)}):\ \rvc\in\mathcal P,\ 1\le j\le n,\ R(\rvw^{(j)})=1\Big\},\qquad \rvw^{(1)},\dots,\rvw^{(n)}\overset{\text{iid}}{\sim}p_\theta(\cdot\mid\rvc),
\label{eq:accept}\\
\theta'&=\text{the point reached by descending }\mathcal L_{\rm SFT}\big(\vartheta;\mathcal A(\theta)\big)\text{ from }\vartheta=\theta,
\label{eq:sft_step}
\end{align}
where $\mathcal A(\theta)$ is first reduced in place, by removing duplicate responses of a prompt and keeping at most $M$ per prompt (the reduced set is the one~(\ref{eq:sft_step}) and Algorithm~\ref{alg:loop} use), the descent runs for a fixed number of epochs, and $\mathcal L_{\rm SFT}$ is ELF's training loss on the accepted pairs, the objective that produced the task-SFT model \citep{hu2026elf}: for $(\rvc,\rvw)\in\mathcal A$ with clean embeddings $\rvx$ of $\rvw$, a random time $t$ and noise $\epsilon$ on the generated positions, $\rvz_t=(1-t)\epsilon+t\rvx$ on $\mathcal G$ with the prompt positions held at $\rvc$,
\begin{equation}
\mathcal L_{\rm SFT}(\vartheta;\mathcal A)=\mathbb E_{(\rvc,\rvw)\in\mathcal A}\Big[\,\mathbb E_{t,\epsilon}\sum_{i\in\mathcal G}\big\|\mathrm d_{\vartheta}(\rvz_t,t,\rvx_{\rm prev},\rvc)[i]-\rvx[i]\big\|^2-\sum_{i\in\mathcal G}\log p_{\vartheta}(w_i\mid\rvx)\Big],
\label{eq:sft}
\end{equation}
the regression of the clean embeddings on the generated positions plus the cross-entropy of the decoder on the clean embeddings. The self-conditioning input $\rvx_{\rm prev}$ is, as in ELF's training, the network's own estimate from a preceding call without gradient for half of the examples and zero for the rest, and the prompt is dropped at random for guidance training; both are kept as in the base training. Rounds are chained: $\theta^{(\varrho+1)}$ is~(\ref{eq:sft_step}) applied to $\theta^{(\varrho)}$ for rounds $\varrho=0,\dots,n_{\rm r}-1$, and the final round's model is the result (Algorithm~\ref{alg:loop}).

\paragraph{What a round does.}
A round is rejection-sampling fine-tuning \citep{yuan2023rft,singh2023restem,zelikman2022star} with the RL policy as the sampler: the accepted responses of a prompt are samples from the policy conditioned on success, and fitting the model to them raises the likelihood of successful responses under the model's own training objective. Two properties of~(\ref{eq:sft}) are the design reasons for using it rather than a sampler-specific loss: it is the objective the model was built with, so a round is ordinary fine-tuning, and it trains the denoiser at all times $t$, which is the mechanism by which the improvement reaches every sampler rather than only the one that produced the samples. A round gives no direct signal on a prompt with no accepted response, so its coverage, the fraction of the pool with at least one success, bounds what the round can teach directly. This is where the critic enters.

\paragraph{Critic-guided sampling.}
The 2D critic is a function of partial states, so it can rank and improve candidates before they are decoded. Both operations use the per-position readings $r_{k,i}$ of~(\ref{eq:value2}) through the score
\begin{equation}
R_{\min}(s_k)=\min_{i\in\mathcal G}r_{k,i},
\label{eq:minscore}
\end{equation}
the most pessimistic of the $|\mathcal G|$ readings of the outcome. The minimum is used here, and the mean in training, because the two tasks differ: the actor improves the policy on average and needs a gradient at every position, whereas a single candidate is accepted only if the whole response is right, so what limits it is its least confident reading, and a gradient of the minimum acts on that reading. \emph{Pruning:} for each prompt, $n_{\rm g}$ trajectories are run in parallel to a step $k_p$; the trajectory with the largest $R_{\min}(s_{k_p})$ continues alone. \emph{Guidance:} at every stochastic step $k=1,\dots,K-1$ (it needs $\gamma>0$), before the denoiser call, the re-noised embeddings are moved along the gradient of the score on the generated positions,
\begin{equation}
\rvz_{t_k}\ \leftarrow\ \rvz_{t_k}+\kappa\,\psi_kc_\sigma\sqrt L\ \frac{\nabla_{\rvz}R_{\min}(s_k)}{\|\nabla_{\rvz}R_{\min}(s_k)\|},
\label{eq:guide}
\end{equation}
a step whose length is proportional to the noise scale $\psi_kc_\sigma$ of that step, a fraction $\kappa\sqrt{L/(|\mathcal G|d)}$ of the norm $\psi_kc_\sigma\sqrt{|\mathcal G|d}$ of the noise the sampler injects there, so that the pushed state stays within the spread of the sampler. The gradient with respect to the state in~(\ref{eq:guide}) is the same kind of quantity as the critic's slope that the actor multiplies with the denoiser Jacobian in~(\ref{eq:pwgrad}); guidance applies it to the sample at inference instead of to the parameters at training. A guided response costs $n_{\rm g}$ trajectories up to step $k_p$, one afterwards, and one backward pass through the critic per step; guided pools are therefore drawn with a cheap sampler (small $K$ and $\gamma$), and the search is meant to compensate the cheap sampler's lower accuracy. Ranking states of a sampler other than the critic's training sampler is an assumption of this use (Appendix~\ref{app:joint_details}). Verification and fine-tuning are unchanged.

\begin{algorithm}[H]
\caption{Amortization loop (with optional critic-guided sampling)}
\label{alg:loop}
\small
\begin{algorithmic}[1]
\REQUIRE policy $\theta^{(0)}$; optional critic $V_\phi$ of \methodSVGTwoD{}; verifier $R$; prompt pool $\mathcal P$; per round: a sampler and whether guidance is used; responses per prompt $n$; cap $M$; rounds $n_{\rm r}$; guidance constants $(n_{\rm g},k_p,\kappa)$
\FOR{round $\varrho=0,\dots,n_{\rm r}-1$}
  \FOR{each prompt $\rvc\in\mathcal P$ and $j=1,\dots,n$}
    \STATE Run the round's sampler under $\theta^{(\varrho)}$; with guidance, apply~(\ref{eq:guide}) at every step and prune $n_{\rm g}\to1$ at step $k_p$ by~(\ref{eq:minscore}); decode $\rvw^{(j)}$
  \ENDFOR
  \STATE $\mathcal A\leftarrow$ the verified, distinct responses, at most $M$ per prompt \hfill (eq.~\ref{eq:accept})
  \STATE $\theta^{(\varrho+1)}\leftarrow$ descend $\mathcal L_{\rm SFT}(\vartheta;\mathcal A)$ from $\vartheta=\theta^{(\varrho)}$ for the round's epochs \hfill (eq.~\ref{eq:sft_step})
\ENDFOR
\STATE \textbf{return} $\theta^{(n_{\rm r})}$
\end{algorithmic}
\end{algorithm}
